\documentclass[10pt,a4paper]{amsart}
\usepackage{amsmath,amssymb,amsthm,graphicx}
% Retained minimal preamble for the Stacks Project chapter more-algebra.tex (More on Algebra).
% The chapter file pinned-original-source.tex carries only \input{preamble} on its first line, so no
% preamble text lives inside the chapter itself.  The Stacks Project uses one single shared preamble for
% every chapter; the snapshot stages it identically as sources/stacks-chow/preamble.tex and
% sources/stacks-etale-cohomology/preamble.tex.  The block below is copied mechanically, line for line,
% from sources/stacks-chow/preamble.tex lines 14-18 (the Xy-pic setup) and 154-217 (the native theorem
% declarations together with the mathematical macros).  Nothing was added, removed or redefined; the
% wrapper class and the amsmath/amssymb/amsthm/graphicx base packages are supplied by the assembler.
% Chapter: more-algebra.tex -- More on Algebra.
% For commutative diagrams we use Xy-pic
\usepackage[all]{xy}

% We use 2cell for 2-commutative diagrams.
\xyoption{2cell}

% Theorem environments.
%
\theoremstyle{plain}
\newtheorem{theorem}[subsection]{Theorem}
\newtheorem{proposition}[subsection]{Proposition}
\newtheorem{lemma}[subsection]{Lemma}

\theoremstyle{definition}
\newtheorem{definition}[subsection]{Definition}
\newtheorem{example}[subsection]{Example}
\newtheorem{exercise}[subsection]{Exercise}
\newtheorem{situation}[subsection]{Situation}

\theoremstyle{remark}
\newtheorem{remark}[subsection]{Remark}
\newtheorem{remarks}[subsection]{Remarks}

\numberwithin{equation}{subsection}

% Macros
%
\def\lim{\mathop{\mathrm{lim}}\nolimits}
\def\colim{\mathop{\mathrm{colim}}\nolimits}
\def\Spec{\mathop{\mathrm{Spec}}}
\def\Hom{\mathop{\mathrm{Hom}}\nolimits}
\def\Ext{\mathop{\mathrm{Ext}}\nolimits}
\def\SheafHom{\mathop{\mathcal{H}\!\mathit{om}}\nolimits}
\def\SheafExt{\mathop{\mathcal{E}\!\mathit{xt}}\nolimits}
\def\Sch{\mathit{Sch}}
\def\Mor{\mathop{\mathrm{Mor}}\nolimits}
\def\Ob{\mathop{\mathrm{Ob}}\nolimits}
\def\Sh{\mathop{\mathit{Sh}}\nolimits}
\def\NL{\mathop{N\!L}\nolimits}
\def\CH{\mathop{\mathrm{CH}}\nolimits}
\def\proetale{{pro\text{-}\acute{e}tale}}
\def\etale{{\acute{e}tale}}
\def\QCoh{\mathit{QCoh}}
\def\Ker{\mathop{\mathrm{Ker}}}
\def\Im{\mathop{\mathrm{Im}}}
\def\Coker{\mathop{\mathrm{Coker}}}
\def\Coim{\mathop{\mathrm{Coim}}}

% Boxtimes
%
\DeclareMathSymbol{\boxtimes}{\mathbin}{AMSa}{"02}

%
% Macros for moduli stacks/spaces
%
\def\QCohstack{\mathcal{QC}\!\mathit{oh}}
\def\Cohstack{\mathcal{C}\!\mathit{oh}}
\def\Spacesstack{\mathcal{S}\!\mathit{paces}}
\def\Quotfunctor{\mathrm{Quot}}
\def\Hilbfunctor{\mathrm{Hilb}}
\def\Curvesstack{\mathcal{C}\!\mathit{urves}}
\def\Polarizedstack{\mathcal{P}\!\mathit{olarized}}
\def\Complexesstack{\mathcal{C}\!\mathit{omplexes}}
% \Pic is the operator that assigns to X its picard group, usage \Pic(X)
% \Picardstack_{X/B} denotes the Picard stack of X over B
% \Picardfunctor_{X/B} denotes the Picard functor of X over B
\def\Pic{\mathop{\mathrm{Pic}}\nolimits}
\def\Picardstack{\mathcal{P}\!\mathit{ic}}
\def\Picardfunctor{\mathrm{Pic}}
\def\Deformationcategory{\mathcal{D}\!\mathit{ef}}

% --- appended 2026-10-12 (Chief): the project own definitions of the reference/slogan/history
% display environments, copied mechanically from the shared Stacks preamble lines 8-12
% (\usepackage{verbatim} plus the three \newenvironment lines, which the chapter preamble keeps
% verbatim from the project preamble). Nothing added, removed or redefined: these environments
% carry no mathematical content and the project itself typesets them as comments. Without these
% three lines a counted statement whose own text uses them fails with Environment slogan undefined.
\usepackage{verbatim}
\newenvironment{reference}{\comment}{\endcomment}
\newenvironment{slogan}{\comment}{\endcomment}
\newenvironment{history}{\comment}{\endcomment}
\usepackage{cleveref}
\crefname{originallemma}{original lemma}{original lemmas}
\Crefname{originallemma}{Original lemma}{Original lemmas}
\makeatletter
\newlabel{more-algebra-proposition-fs-flat-fibre-fs}{{07NQ}{source}{}{sourcepointer-1}{}}
\newlabel{more-algebra-proposition-fs-flat-fibre-fs@cref}{{[originallemma][0][]07NQ}{[1][0][]source}{}{}{}}
\makeatother
\crefname{originallemma}{original lemma}{original lemmas}
\Crefname{originallemma}{Original lemma}{Original lemmas}
\makeatletter
\newlabel{theorem-regular-fs}{{07EL}{source}{}{sourcepointer-2}{}}
\newlabel{theorem-regular-fs@cref}{{[originallemma][0][]07EL}{[1][0][]source}{}{}{}}
\makeatother
\crefname{originallemma}{original lemma}{original lemmas}
\Crefname{originallemma}{Original lemma}{Original lemmas}
\makeatletter
\newlabel{lemma-formally-smooth-flat}{{07NP}{source}{}{sourcepointer-3}{}}
\newlabel{lemma-formally-smooth-flat@cref}{{[originallemma][0][]07NP}{[1][0][]source}{}{}{}}
\makeatother
\crefname{originallemma}{original lemma}{original lemmas}
\Crefname{originallemma}{Original lemma}{Original lemmas}
\makeatletter
\newlabel{lemma-base-change-fs}{{07EG}{source}{}{sourcepointer-4}{}}
\newlabel{lemma-base-change-fs@cref}{{[originallemma][0][]07EG}{[1][0][]source}{}{}{}}
\makeatother
\crefname{originallemma}{original lemma}{original lemmas}
\Crefname{originallemma}{Original lemma}{Original lemmas}
\makeatletter
\newlabel{lemma-formally-smooth-completion}{{07ED}{source}{}{sourcepointer-5}{}}
\newlabel{lemma-formally-smooth-completion@cref}{{[originallemma][0][]07ED}{[1][0][]source}{}{}{}}
\makeatother
\crefname{originallemma}{original lemma}{original lemmas}
\Crefname{originallemma}{Original lemma}{Original lemmas}
\makeatletter
\newlabel{algebra-lemma-completion-faithfully-flat}{{00MC}{source}{}{sourcepointer-6}{}}
\newlabel{algebra-lemma-completion-faithfully-flat@cref}{{[originallemma][0][]00MC}{[1][0][]source}{}{}{}}
\makeatother
\crefname{originallemma}{original lemma}{original lemmas}
\Crefname{originallemma}{Original lemma}{Original lemmas}
\makeatletter
\newlabel{lemma-embed-map-Noetherian-complete-local-rings}{{07NN}{source}{}{sourcepointer-7}{}}
\newlabel{lemma-embed-map-Noetherian-complete-local-rings@cref}{{[originallemma][0][]07NN}{[1][0][]source}{}{}{}}
\makeatother
\crefname{originallemma}{original lemma}{original lemmas}
\Crefname{originallemma}{Original lemma}{Original lemmas}
\makeatletter
\newlabel{algebra-lemma-regular-quotient-regular}{{00NR}{source}{}{sourcepointer-8}{}}
\newlabel{algebra-lemma-regular-quotient-regular@cref}{{[originallemma][0][]00NR}{[1][0][]source}{}{}{}}
\makeatother
\crefname{originallemma}{original lemma}{original lemmas}
\Crefname{originallemma}{Original lemma}{Original lemmas}
\makeatletter
\newlabel{lemma-quotient-power-series-ring-over-Cohen}{{07NM}{source}{}{sourcepointer-9}{}}
\newlabel{lemma-quotient-power-series-ring-over-Cohen@cref}{{[originallemma][0][]07NM}{[1][0][]source}{}{}{}}
\makeatother
\crefname{originallemma}{original lemma}{original lemmas}
\Crefname{originallemma}{Original lemma}{Original lemmas}
\makeatletter
\newlabel{algebra-lemma-mod-injective}{{00ME}{source}{}{sourcepointer-10}{}}
\newlabel{algebra-lemma-mod-injective@cref}{{[originallemma][0][]00ME}{[1][0][]source}{}{}{}}
\makeatother
\crefname{originallemma}{original lemma}{original lemmas}
\Crefname{originallemma}{Original lemma}{Original lemmas}
\makeatletter
\newlabel{definition-formally-smooth}{{07EB}{source}{}{sourcepointer-11}{}}
\newlabel{definition-formally-smooth@cref}{{[originallemma][0][]07EB}{[1][0][]source}{}{}{}}
\makeatother
\crefname{originallemma}{original lemma}{original lemmas}
\Crefname{originallemma}{Original lemma}{Original lemmas}
\makeatletter
\newlabel{lemma-power-series-ring-over-Cohen-fs}{{07NL}{source}{}{sourcepointer-12}{}}
\newlabel{lemma-power-series-ring-over-Cohen-fs@cref}{{[originallemma][0][]07NL}{[1][0][]source}{}{}{}}
\makeatother
\crefname{originallemma}{original lemma}{original lemmas}
\Crefname{originallemma}{Original lemma}{Original lemmas}
\makeatletter
\newlabel{algebra-lemma-universal-omega}{{00RO}{source}{}{sourcepointer-13}{}}
\newlabel{algebra-lemma-universal-omega@cref}{{[originallemma][0][]00RO}{[1][0][]source}{}{}{}}
\makeatother
\crefname{originallemma}{original lemma}{original lemmas}
\Crefname{originallemma}{Original lemma}{Original lemmas}
\makeatletter
\newlabel{algebra-proposition-characterize-separable-field-extensions}{{0322}{source}{}{sourcepointer-14}{}}
\newlabel{algebra-proposition-characterize-separable-field-extensions@cref}{{[originallemma][0][]0322}{[1][0][]source}{}{}{}}
\makeatother
\title{Flatness and formal smoothness of a local map versus regularity of its special fibre}
\author{The Stacks Project Authors}
\date{stacks-project:more-algebra.tex:local-snapshot-20261012; GNU FDL 1.2 or later}
\begin{document}
\maketitle
\setcounter{subsection}{1}
\noindent\textit{Editorial note.} All delivered mathematical text is the author's own, copied from the pinned Stacks Project chapter file more-algebra.tex; no statement or proof step was rewritten, reordered or removed.\par\medskip
\noindent\textit{Editorial note.} The chapter file pinned-original-source.tex carries only \textbackslash{}input\{preamble\} on its first line, so no preamble text lives inside the chapter. The retained file sources/stacks-more-algebra/retained-preamble.tex is a mechanical, line-for-line copy of the project's single shared preamble as staged in the snapshot (Xy-pic setup plus theorem declarations and mathematical macros). Nothing was added, removed or redefined, and the excerpt serves only as preamble, never as a span coordinate.\par\medskip
\noindent\textit{Editorial note.} Both counted spans start at the \textbackslash{}begin line of their environment, so the statement and the author's proof are delivered whole; surrounding prose of the section is deliberately outside the counted spans and is carried as standard prerequisites instead.\par\medskip
\noindent\textit{Editorial note.} The citation fields (source\_id, source\_version, authors, source\_url, license.name, source\_title) are identical to those of every other card this lane has produced from this chapter, so the batch selector treats them as one source.\par\medskip
\noindent\textit{Allowed standard prerequisites (editorial).}\begin{itemize}
\item Equivalence of geometric regularity of the fibre, formal smoothness of the fibre, and formal smoothness of the local map. Theorem \textbackslash{}ref\{theorem-regular-fs\}, cited at proof lines 10096 and again at 10198 and 10241. Gives the equivalence of (1) and (2), and supplies the final lifting argument used to finish (2) implies (3).
\item Formal smoothness of a local homomorphism implies flatness. Lemma \textbackslash{}ref\{lemma-formally-smooth-flat\}, cited at proof line 10099. Yields flatness of A to B from condition (3).
\item Base change of formal smoothness, and preservation of formal smoothness under completion. Lemma \textbackslash{}ref\{lemma-base-change-fs\} and Lemma \textbackslash{}ref\{lemma-formally-smooth-completion\}, cited at proof lines 10100, 10104, and 10150. Transport the hypothesis to the special fibre and let the author replace A and B by their completions.
\item Existence of the diagram of Noetherian complete local rings S over R mapping to B over A. Lemma \textbackslash{}ref\{lemma-embed-map-Noetherian-complete-local-rings\}, cited at proof lines 10112 and 10143. Provides the power series model of the extension that the reduction to equal special fibres works with.
\item The quotient of a power series ring by part of a regular system of parameters is again a power series ring over a field or a Cohen ring. Lemma \textbackslash{}ref\{lemma-quotient-power-series-ring-over-Cohen\}, cited at proof lines 10135 and 10165. Identifies S after the cut as a power series ring, which is what makes the universal property argument available.
\end{itemize}\medskip
\ifdefined\hypertarget\hypertarget{sourcepointer-1}{}\fi
\noindent\textit{Source reference.} Original lemma 07NQ: focused-seeds/AGP/sources/stacks-snapshot-20261011/tags.txt, source lines 9770--9770. The lemma's own canonical tag, quoted so that readers can resolve the label inside the delivered text; it points at the counted statement itself and carries no proof content.\par\medskip
\ifdefined\hypertarget\hypertarget{sourcepointer-2}{}\fi
\noindent\textit{Source reference.} Original lemma 07EL: focused-seeds/AGP/sources/stacks-snapshot-20261011/tags.txt, source lines 9451--9451. Navigation pointer to a statement the counted proof refers to but does not contain; nothing load-bearing is delegated to it.\par\medskip
\ifdefined\hypertarget\hypertarget{sourcepointer-3}{}\fi
\noindent\textit{Source reference.} Original lemma 07NP: focused-seeds/AGP/sources/stacks-snapshot-20261011/tags.txt, source lines 9769--9769. Navigation pointer to a statement the counted proof refers to but does not contain; nothing load-bearing is delegated to it.\par\medskip
\ifdefined\hypertarget\hypertarget{sourcepointer-4}{}\fi
\noindent\textit{Source reference.} Original lemma 07EG: focused-seeds/AGP/sources/stacks-snapshot-20261011/tags.txt, source lines 9446--9446. Navigation pointer to a statement the counted proof refers to but does not contain; nothing load-bearing is delegated to it.\par\medskip
\ifdefined\hypertarget\hypertarget{sourcepointer-5}{}\fi
\noindent\textit{Source reference.} Original lemma 07ED: focused-seeds/AGP/sources/stacks-snapshot-20261011/tags.txt, source lines 9443--9443. Navigation pointer to a statement the counted proof refers to but does not contain; nothing load-bearing is delegated to it.\par\medskip
\ifdefined\hypertarget\hypertarget{sourcepointer-6}{}\fi
\noindent\textit{Source reference.} Original lemma 00MC: focused-seeds/AGP/sources/stacks-snapshot-20261011/tags.txt, source lines 818--818. Navigation pointer to a statement the counted proof refers to but does not contain; nothing load-bearing is delegated to it.\par\medskip
\ifdefined\hypertarget\hypertarget{sourcepointer-7}{}\fi
\noindent\textit{Source reference.} Original lemma 07NN: focused-seeds/AGP/sources/stacks-snapshot-20261011/tags.txt, source lines 9768--9768. Navigation pointer to a statement the counted proof refers to but does not contain; nothing load-bearing is delegated to it.\par\medskip
\ifdefined\hypertarget\hypertarget{sourcepointer-8}{}\fi
\noindent\textit{Source reference.} Original lemma 00NR: focused-seeds/AGP/sources/stacks-snapshot-20261011/tags.txt, source lines 869--869. Navigation pointer to a statement the counted proof refers to but does not contain; nothing load-bearing is delegated to it.\par\medskip
\ifdefined\hypertarget\hypertarget{sourcepointer-9}{}\fi
\noindent\textit{Source reference.} Original lemma 07NM: focused-seeds/AGP/sources/stacks-snapshot-20261011/tags.txt, source lines 9767--9767. Navigation pointer to a statement the counted proof refers to but does not contain; nothing load-bearing is delegated to it.\par\medskip
\ifdefined\hypertarget\hypertarget{sourcepointer-10}{}\fi
\noindent\textit{Source reference.} Original lemma 00ME: focused-seeds/AGP/sources/stacks-snapshot-20261011/tags.txt, source lines 820--820. Navigation pointer to a statement the counted proof refers to but does not contain; nothing load-bearing is delegated to it.\par\medskip
\ifdefined\hypertarget\hypertarget{sourcepointer-11}{}\fi
\noindent\textit{Source reference.} Original lemma 07EB: focused-seeds/AGP/sources/stacks-snapshot-20261011/tags.txt, source lines 9441--9441. Navigation pointer to a statement the counted proof refers to but does not contain; nothing load-bearing is delegated to it.\par\medskip
\ifdefined\hypertarget\hypertarget{sourcepointer-12}{}\fi
\noindent\textit{Source reference.} Original lemma 07NL: focused-seeds/AGP/sources/stacks-snapshot-20261011/tags.txt, source lines 9766--9766. Navigation pointer to a statement the counted proof refers to but does not contain; nothing load-bearing is delegated to it.\par\medskip
\ifdefined\hypertarget\hypertarget{sourcepointer-13}{}\fi
\noindent\textit{Source reference.} Original lemma 00RO: focused-seeds/AGP/sources/stacks-snapshot-20261011/tags.txt, source lines 1010--1010. Navigation pointer to a statement the counted proof refers to but does not contain; nothing load-bearing is delegated to it.\par\medskip
\ifdefined\hypertarget\hypertarget{sourcepointer-14}{}\fi
\noindent\textit{Source reference.} Original lemma 0322: focused-seeds/AGP/sources/stacks-snapshot-20261011/tags.txt, source lines 4005--4005. Navigation pointer to a statement the counted proof refers to but does not contain; nothing load-bearing is delegated to it.\par\medskip
\begin{proposition}
\label{proposition-fs-flat-fibre-fs}
Let $A \to B$ be a local homomorphism of Noetherian local rings.
Let $k$ be the residue field of $A$ and $\overline{B} = B \otimes_A k$
the special fibre. The following are equivalent
\begin{enumerate}
\item $A \to B$ is flat and $\overline{B}$ is geometrically regular
over $k$,
\item $A \to B$ is flat and $k \to \overline{B}$ is formally smooth
in the $\mathfrak m_{\overline{B}}$-adic topology, and
\item $A \to B$ is formally smooth in the $\mathfrak m_B$-adic
topology.
\end{enumerate}
\end{proposition}

\begin{proof}
The equivalence of (1) and (2) follows from Theorem \ref{theorem-regular-fs}.

\medskip\noindent
Assume (3). By Lemma \ref{lemma-formally-smooth-flat} we see that
$A \to B$ is flat. By Lemma \ref{lemma-base-change-fs} we see that
$k \to \overline{B}$ is formally smooth in the
$\mathfrak m_{\overline{B}}$-adic topology. Thus (2) holds.

\medskip\noindent
Assume (2). Lemma \ref{lemma-formally-smooth-completion}
tells us formal smoothness is preserved under completion. The same is true
for flatness by Algebra, Lemma \ref{algebra-lemma-completion-faithfully-flat}.
Hence we may replace $A$ and $B$ by their respective completions and
assume that $A$ and $B$ are Noetherian complete local rings.
In this case choose a diagram
$$
\xymatrix{
S \ar[r] & B \\
R \ar[u] \ar[r] & A \ar[u]
}
$$
as in Lemma \ref{lemma-embed-map-Noetherian-complete-local-rings}.
We will use all of the properties of this diagram without further mention.
Fix a regular system of parameters $t_1, \ldots, t_d$ of $R$
with $t_1 = p$ in case the characteristic of $k$ is $p > 0$.
Set $\overline{S} = S \otimes_R k$. Consider the short exact sequence
$$
0 \to J \to S \to B \to 0
$$
As $\overline{B}$ and $\overline{S}$ are regular, the kernel of
$\overline{S} \to \overline{B}$ is generated by elements
$\overline{x}_1, \ldots, \overline{x}_r$ which form part of a regular
system of parameters of $\overline{S}$, see
Algebra, Lemma \ref{algebra-lemma-regular-quotient-regular}.
Lift these elements to $x_1, \ldots, x_r \in J$. Then
$t_1, \ldots, t_d, x_1, \ldots, x_r$ is part of a regular system of
parameters for $S$. Hence $S/(x_1, \ldots, x_r)$ is a power
series ring over a field (if the characteristic of $k$ is zero)
or a power series ring over a Cohen ring (if the characteristic of
$k$ is $p > 0$), see
Lemma \ref{lemma-quotient-power-series-ring-over-Cohen}.
Moreover, it is still the case that $R \to S/(x_1, \ldots, x_r)$
maps $t_1, \ldots, t_d$ to a part of a regular system of parameters
of $S/(x_1, \ldots, x_r)$. In other words, we may replace $S$ by
$S/(x_1, \ldots, x_r)$ and assume we have a diagram
$$
\xymatrix{
S \ar[r] & B \\
R \ar[u] \ar[r] & A \ar[u]
}
$$
as in Lemma \ref{lemma-embed-map-Noetherian-complete-local-rings}
with moreover $\overline{S} = \overline{B}$. In this case the map
$$
S \otimes_R A \longrightarrow B
$$
is an isomorphism as it is surjective, an isomorphism on special
fibres, and source and target are flat over $A$ (for example use
Algebra, Lemma \ref{algebra-lemma-mod-injective} or use that
tensoring the short exact sequence $0 \to I \to S \otimes_R A \to B \to 0$
over $A$ with $k$ we find $I \otimes_A k = 0$ hence $I = 0$ by Nakayama).
Thus by Lemma \ref{lemma-base-change-fs}
it suffices to show that $R \to S$ is formally
smooth in the $\mathfrak m_S$-adic topology.
Of course, since $\overline{S} = \overline{B}$, we have
that $\overline{S}$ is formally smooth over $k = R/\mathfrak m_R$.

\medskip\noindent
Choose elements $y_1, \ldots, y_m \in S$ such that
$t_1, \ldots, t_d, y_1, \ldots, y_m$ is a regular system of parameters
for $S$. If the characteristic of $k$ is zero, choose a coefficient
field $K \subset S$ and if the characteristic of $k$ is $p > 0$
choose a Cohen ring $\Lambda \subset S$ with residue field $K$.
At this point the map $K[[t_1, \ldots, t_d, y_1, \ldots, y_m]] \to S$
(characteristic zero case) or
$\Lambda[[t_2, \ldots, t_d, y_1, \ldots, y_m]] \to S$
(characteristic $p > 0$ case) is an isomorphism, see
Lemma \ref{lemma-quotient-power-series-ring-over-Cohen}.
From now on we think of $S$ as the above power series ring.

\medskip\noindent
The rest of the proof is analogous to the argument
in the proof of Theorem \ref{theorem-regular-fs}. Choose a solid diagram
$$
\xymatrix{
S \ar[r]_{\bar\psi} \ar@{-->}[rd] & N/J \\
R \ar[u]^i \ar[r]^\varphi & N \ar[u]_\pi
}
$$
as in Definition \ref{definition-formally-smooth}. As $J^2 = 0$ we see
that $J$ has a canonical $N/J$ module structure and via $\bar\psi$ a
$S$-module structure. As $\bar\psi$ is continuous for the
$\mathfrak m_S$-adic topology we see that $\mathfrak m_S^nJ = 0$ for some
$n$. Hence we can filter $J$ by $N/J$-submodules
$0 \subset J_1 \subset J_2 \subset \ldots \subset J_n = J$
such that each quotient $J_{t + 1}/J_t$ is annihilated by $\mathfrak m_S$.
Considering the sequence of ring maps
$N \to N/J_1 \to N/J_2 \to \ldots \to N/J$
we see that it suffices to prove the existence of the dotted arrow when
$J$ is annihilated by $\mathfrak m_S$, i.e., when $J$ is a
$K$-vector space.

\medskip\noindent
Assume given a diagram as above such that $J$ is annihilated by
$\mathfrak m_S$. As $\mathbf{Q} \to S$ (characteristic zero case)
or $\mathbf{Z} \to S$ (characteristic $p > 0$ case)
is formally smooth in the $\mathfrak m_S$-adic topology (see
Lemma \ref{lemma-power-series-ring-over-Cohen-fs}), we can find
a ring map $\psi : S \to N$ such that $\pi \circ \psi = \bar \psi$.
Since $S$ is a power series ring in $t_1, \ldots, t_d$ (characteristic zero)
or $t_2, \ldots, t_d$ (characteristic $p > 0$) over
a subring, it follows from the universal property of power series rings
that we can change our choice of $\psi$ so that $\psi(t_i)$ equals
$\varphi(t_i)$ (automatic for $t_1 = p$ in the characteristic $p$ case).
Then $\psi \circ i$ and $\varphi : R \to N$ are two maps whose
compositions with $\pi$ are equal and which agree on $t_1, \ldots, t_d$.
Hence $D = \psi \circ i - \varphi : R \to J$ is a derivation which
annihilates $t_1, \ldots, t_d$.
By Algebra, Lemma \ref{algebra-lemma-universal-omega} we can write
$D = \xi \circ \text{d}$ for some $R$-linear map
$\xi : \Omega_{R/\mathbf{Z}} \to J$ which annihilates
$\text{d}t_1, \ldots, \text{d}t_d$ (by construction) and
$\mathfrak m_R \Omega_{R/\mathbf{Z}}$ (as $J$ is annihilated by
$\mathfrak m_R$). Hence $\xi$ factors as a composition
$$
\Omega_{R/\mathbf{Z}} \to \Omega_{k/\mathbf{Z}} \xrightarrow{\xi'} J
$$
where $\xi'$ is $k$-linear. Using the $K$-vector space structure on $J$ we
extend $\xi'$ to a $K$-linear map
$$
\xi'' : \Omega_{k/\mathbf{Z}} \otimes_k K \longrightarrow J.
$$
Using that $\overline{S}/k$ is formally smooth we see that
$$
\Omega_{k/\mathbf{Z}} \otimes_k K \to
\Omega_{\overline{S}/\mathbf{Z}} \otimes_S K
$$
is injective by Theorem \ref{theorem-regular-fs} (this is true also
in the characteristic zero case as it is even true that
$\Omega_{k/\mathbf{Z}} \to \Omega_{K/\mathbf{Z}}$ is injective
in characteristic zero, see Algebra,
Proposition \ref{algebra-proposition-characterize-separable-field-extensions}).
Hence we can find a $K$-linear map
$\xi''' : \Omega_{\overline{S}/\mathbf{Z}} \otimes_S K \to J$ whose
restriction to $\Omega_{k/\mathbf{Z}} \otimes_k K$ is $\xi''$. Write
$$
D' : S \xrightarrow{\text{d}} \Omega_{S/\mathbf{Z}}
\to \Omega_{\overline{S}/\mathbf{Z}} \to
\Omega_{\overline{S}/\mathbf{Z}} \otimes_S K \xrightarrow{\xi'''} J.
$$
Finally, set $\psi' = \psi - D' : S \to N$. The reader verifies that $\psi'$
is a ring map such that $\pi \circ \psi' = \bar \psi$ and such that
$\psi' \circ i = \varphi$ as desired.
\end{proof}

\end{document}
